\subsection{Analytical Assessment and Identifiability}
\label{sec:analytical_assessment}

The framework distinguishes a description of the path from a decision about
performance. We formalize this distinction using a condition vector
$z=(C,R_*,\{R_i\},Q,L,v,A,n,W)$: bottleneck capacity, reference and
per-flow propagation RTTs, absolute queue limit, loss process, implementation,
queue policy, flow count and observation window. The physical descriptors are
$\mathrm{BDP}_*=CR_*$ and $q=Q/(CR_*)$, with capacity and queue expressed
in consistent bit units. These descriptors organize cases; the following
analysis shows why they cannot determine bandwidth allocation by themselves.

\paragraph{Allocation is not identified by the in-flight identity.}
For two flows, let their aggregate rate be $0<G\leq C$ and let $s$ be the
second flow's share. The rates $x_1=(1-s)G$ and $x_2=sG$ satisfy the same
aggregate-rate constraint for every $s\in[0,1]$. Their propagation-pipe
amounts $I_i=x_iR_i$ also satisfy the per-flow rate--delay identities.
Thus, even with $C,R_1,R_2,Q$ fixed, the identities admit different shares;
a longer propagation RTT does not algebraically force a larger share.
This is an identifiability result about the equations, not a claim that all
shares are equilibria of a particular controller. Selecting the realized
allocation requires controller and queue dynamics or direct measurements.

\paragraph{Utilization and fairness are separate constraints.}
Define normalized receiver goodput $u=G/C$ and imbalance
$d=|2s-1|$. Substituting the two rates into Jain's standard index
\citep{jain1984fairness} gives
\begin{equation}
 J=\frac{1}{2\{s^2+(1-s)^2\}}=\frac{1}{1+d^2}.
 \label{eq:share_jain}
\end{equation}
Consequently, a fairness requirement $J\geq j_0$, where
$1/2\leq j_0\leq1$, is equivalent to
\begin{equation}
 \frac{1-\sqrt{j_0^{-1}-1}}{2}\leq s\leq
 \frac{1+\sqrt{j_0^{-1}-1}}{2}.
 \label{eq:share_acceptance}
\end{equation}
For example, $j_0=0.95$ requires a share between 0.385 and 0.615,
regardless of whether aggregate goodput fills the link. Jain's index and
this algebraic transformation are established metrics, not new fairness
measures. Our analytical contribution is to use their explicit feasible
region to separate path descriptors, performance requirements and observable
allocation in a BBR assessment.

\paragraph{A requirement-dependent assessment procedure.}
A workload specifies its requirements before protocol selection: a minimum
$u_0$, a minimum $j_0$, a delay limit $D_0$, and a retransmission budget
$H_0$. For each required metric, the assessment returns \emph{meets},
\emph{does not meet}, or \emph{undetermined}. A missing or incompatible
measurement window yields \emph{undetermined}, rather than an inferred pass.
A joint assessment meets the requirements only if every required metric
meets its own limit. An observed violation yields a failure for that
requirement; if there is no observed violation but some required evidence is
missing, the joint result remains undetermined. A zero-goodput run has $u=0$ and undefined Jain fairness; it cannot pass
a positive utilization requirement. This procedure allows
throughput and fairness evidence to be usable without claiming that they
establish latency or retransmission suitability.

For repeated runs, we report both metric means and the empirical fraction
satisfying the joint rate/share conditions:
\begin{equation}
 \widehat p_{u,J}=\frac{1}{N}\sum_{r=1}^{N}
 \mathbf{1}\{u_r\geq u_0,\ J_r\geq j_0\}.
 \label{eq:assessment_fraction}
\end{equation}
The fraction is descriptive with five repetitions per condition; it is not
a population reliability guarantee. Equation~\ref{eq:share_jain} is applied
within each run. In general, substituting an averaged share does not recover
the mean Jain index. Queue bins and BDP bins do not set $u_0,j_0,D_0,H_0$:
those limits express service requirements and must be reported separately.
Section~\ref{sec:framework_transfer} evaluates the rate/fairness component
and its transfer limitations using the archived measurements.
